\subsection{有限表示模}

A-模 E 的一个表示（或长度为 1 的表示）是一个正合序列

\[
\mathrm{L} _ {1} \rightarrow \mathrm{L} _ {0} \rightarrow \mathrm{E} \rightarrow 0\tag{13}
\]

其中 \(\mathbf{L}_0\) 和 \(\mathbf{L}_1\) 是自由 A-模。

每个 A-模 E 都有一个表示。事实上，我们知道（II，第 27 页，命题 20）存在一个满同态 \(u: \mathbf{L}_0 \to \mathbf{E}\) ，其中 \(\mathbf{L}_0\) 是自由的；如果 R 是 \(u\) 的核，同样存在一个满同态 \(v: \mathbf{L}_1 \to \mathbf{R}\) ，其中 \(\mathbf{L}_1\) 是自由的。如果将 \(v\) 视为从 \(\mathbf{L}_1\) 到 \(\mathbf{L}_0\) 的同态，序列 \(\mathbf{L}_1 \xrightarrow{v} \mathbf{L}_0 \xrightarrow{u} \mathbf{E} \longrightarrow 0\) 依定义是正合的，因此我们的断言成立。

若 \(\rho: A \to B\) 是一个环同态，那么 \(E\) 的每个表示 (13) 给出 \(E_{(B)} = B \otimes_A E\) 的一个表示：

\[
\mathbf {B} \otimes_ {\mathbf {A}} \mathrm{L} _ {1} \rightarrow \mathbf {B} \otimes_ {\mathbf {A}} \mathrm{L} _ {0} \rightarrow \mathbf {B} \otimes_ {\mathbf {A}} \mathrm{E} \rightarrow 0\tag{14}
\]

这由 II，第 58 页，命题 5 以及 \(\mathbf{B} \otimes_{\mathbf{A}} \mathbf{L}\) 当 L 是自由模时是自由 B-模这一事实得出。

称模 E 的一个表示 (13) 是有限的，如果自由模 \(L_{0}\) 和 \(L_{1}\) 有有限基。显然，如果表示 (13) 是有限的，那么表示 (14) 也是有限的。称 E 为有限表示的 A-模，如果它有一个有限表示。

命题5. --- (i) 每个具有有限表示的模都是有限生成的。

\begin{enumerate}
\def\labelenumi{(\roman{enumi})}
\setcounter{enumi}{1}
\item
  若A是左诺特环，则每个有限生成的A-模都有有限表示。
\item
  每个有限生成的投射模都具有有限表示。
\end{enumerate}

断言（i）由定义直接得出。假设A是左诺特的，且E是有限生成的。那么存在一个满同态 \(u: L_0 \to E\) ，其中 \(L_0\) 是具有有限基的自由 A-模；\(u\) 的核 R 是有限生成的，因此存在一个满同态 \(v: L_1 \to R\) ，其中 \(L_1\) 是具有有限基的自由模，并且正合序列 \(L_1 \xrightarrow{v} L_0 \xrightarrow{u} E \longrightarrow 0\) 是E的一个有限表示；由此得(ii)。

最后，假设 E 是一个有限生成的投射模；那么它是有限生成自由模 \(L_{0}\) 的一个直和项（II，第 40 页，推论 1）；满同态 \(L_{0} \rightarrow E\) 的核 R 则同构于 \(L_{0}\) 的一个商，因此是有限生成的，并像上面那样完成证明。

命题6. --- 设A是一个环，E是一个有限表示的A-模。对于每个正合序列

\[
0 \longrightarrow \mathrm{F} \stackrel {j} {\longrightarrow} \mathrm{G} \stackrel {p} {\longrightarrow} \mathrm{E} \longrightarrow 0
\]

其中G是有限生成的，模F也是有限生成的。

设 \(L_{1} \xrightarrow{r} L_{0} \xrightarrow{s} E \longrightarrow 0\) 是一个有限表示；如果 \((e_{i})\) 是 \(L_{0}\) 的一组基，则对每个 i 存在一个元素 \(g_{i} \in G\) 使得 \(p(g_{i}) = s(e_{i})\) ；同态 \(u : L_{0} \to G\) （对一切 i 满足 \(u(e_{i}) = g_{i}\) ）因而满足 \(s = p \circ u\) 。由于 \(s \circ r = 0\) ，我们有 \(u(r(L_{1})) \subset \operatorname{Ker} p\) ，并且由于 Ker p 同构于 F，我们看到存在一个同态 \(v : L_{1} \to F\) 使得图表

\[
\begin{array}{c c c c c c c} L _ {1} & \xrightarrow {r} & L _ {0} & \xrightarrow {s} & E & \longrightarrow & 0 \\ v \Big \downarrow & & u \Big \downarrow & & ^ {1 _ {E}} \Big \downarrow \\ F & \xrightarrow {} _ {j} & G & \xrightarrow {} _ {p} & E & \longrightarrow & 0 \end{array}
\]

是交换的。由于 j 是单射且 s 是满射，可以应用 X 第 4 页的命题 2，换言之存在一个正合序列

\[
\operatorname{Ker} 1 _ {\mathrm{E}} \xrightarrow {d} \text { Coker } v \longrightarrow \text { Coker } u \longrightarrow \text { Coker } 1 _ {\mathrm{E}}.
\]

这表明 Coker v 同构于 \(\mathrm{G}/u(\mathrm{L}_{0})\) ，而根据假设它是有限生成的。此外，我们还有正合序列

\[
0 \rightarrow v (\mathbf {L} _ {1}) \rightarrow \mathrm{F} \rightarrow \text { Coker } v \rightarrow 0
\]

并且由于 \(v(\mathbf{L}_1)\) 和Coker \(v\) 是有限生成的，所以 F（II，第 17 页，推论 5）也是有限生成的。

命题 7. —— 设 M 是一个 A-模。存在一个右有向有序集 I 和一个相对于 I 的、由有限表示 A-模构成的归纳系（M \(_{\alpha}\) , φ \(_{\beta\alpha}\) ），使得 M 同构于 lim M \(_{\alpha}\) 。若 M 具有由 n 个元素组成的生成族，则可假定 M \(_{\alpha}\)

亦然。考虑一个表示

\[
\mathrm{A} _ {s} ^ {(\mathrm{K})} \xrightarrow {u} \mathrm{A} _ {s} ^ {(\mathrm{L})} \xrightarrow {v} \mathrm{M} \longrightarrow 0;
\]

设 I 为由对 \(\alpha = (\mathbf{K}', \mathbf{L}')\) 组成的集合，其中 \(\mathbf{K}'\) （相应地， \(\mathbf{L}'\) ）是 \(\mathbf{K}\) （相应地 L）的有限子集，使得 \(u\) 诱导出映射 \(u_{\alpha}\) ，它把子模 \(\mathbf{A}_s^{\mathbf{K}'}\) （\(\mathbf{A}_s^{(\mathbf{K})}\) 的子模）映入子模 \(\mathbf{A}_s^{\mathbf{L}'}\) （\(\mathbf{A}_s^{(\mathbf{L})}\) 的子模）；对于 \(\alpha \in \mathbf{I}\) ，设 \(\mathbf{M}_{\alpha}\) 为 \(u_{\alpha}\) 的余核，并设 \(v_{\alpha}: \mathbf{A}_s^{\mathbf{L}'} \to \mathbf{M}_{\alpha}\) 为典范映射，从而我们有一个行正合的交换图：

\[
\mathrm{A} _ {s} ^ {(\mathrm{K})} \xrightarrow {u} \mathrm{A} _ {s} ^ {(\mathrm{L})} \xrightarrow {i} \mathrm{M} \longrightarrow 0
\]

\[
i _ {\alpha} \uparrow , \quad j _ {\alpha} \uparrow , \quad f _ {\alpha} \uparrow
\]

\[
\mathrm{A} _ {s} ^ {\mathrm{K} ^ {\prime}} \xrightarrow {u _ {\alpha}} \mathrm{A} _ {s} ^ {\mathrm{L} ^ {\prime}} \xrightarrow {v _ {\alpha}} \mathrm{M} _ {\alpha} \longrightarrow 0,
\]

其中 \(i_{\alpha}\) 和 \(j_{\alpha}\) 是典范单射，而 \(f_{\alpha}\) 是由 \(j_{\alpha}\) 通过取商得出的。按以下关系给集合 I 排序：

\[
\alpha = \left(\mathrm{K} ^ {\prime}, \mathrm{L} ^ {\prime}\right) \leqslant \beta = \left(\mathrm{K} ^ {\prime \prime}, \mathrm{L} ^ {\prime \prime}\right) \quad \text { if } \quad \mathrm{K} ^ {\prime} \subset \mathrm{K} ^ {\prime \prime}, \quad \mathrm{L} ^ {\prime} \subset \mathrm{L} ^ {\prime \prime};
\]

当 \(\alpha \leqslant \beta\) ，设 \(\varphi_{\beta \alpha}: \mathbf{M}_{\alpha} \to \mathbf{M}_{\beta}\) 为由 \(\mathbf{A}_s^{L'}\) 到 \(\mathbf{A}_s^{L''}\) 的包含映射通过取商得出的同态。于是可直接验证有序集 I 是有向的，且 \((\mathbf{M}_{\alpha}, \varphi_{\beta \alpha})\) 是 A-模的归纳系， \((\varphi_{\alpha})\) 是 A-同态的归纳系。取归纳极限，得到一个交换图

\[
\begin{array}{c c c c c} \mathbf {A} _ {s} ^ {(\mathbf {K})} & \stackrel {{u}} {{\longrightarrow}} & \mathbf {A} _ {s} ^ {(L)} & \stackrel {{v}} {{\longrightarrow}} & \mathbf {M} \\ i \uparrow & & j \uparrow & & \varphi \uparrow \\ \underline {{\lim}} \mathbf {A} _ {s} ^ {\mathbf {K} ^ {\prime}} & \longrightarrow & \underline {{\lim}} \mathbf {A} _ {s} ^ {\mathbf {L} ^ {\prime}} & \longrightarrow & \underline {{\lim}} \mathbf {M} _ {\alpha} \longrightarrow 0; \end{array}\tag{15}
\]

此图的行是正合的（II，p.～91，命题 3）；由于 i 和 j 是双射，φ 也是双射（X，p.～7，推论 3），由此得命题。

